Klinische Leitlinien geben Empfehlungen für die Diagnose, Behandlung und Überwachung von Patientinnen und Patienten, doch verschiedene Organisationen stellen ähnliche Informationen häufig auf unterschiedliche Weise dar. Sie verwenden unter Umständen abweichende Terminologien, Dokumentstrukturen, Empfehlungsformate und Maßeinheiten. In der Diabetesversorgung kann dies es erschweren, Informationen aus mehreren Leitlinien zusammenzuführen und sie einheitlich für Aufgaben wie die Suche nach klinischen Informationen oder die Generierung von Lernfragen zu nutzen.

Diese Arbeit untersucht, wie Diabetesleitlinien verschiedener Organisationen in einem gemeinsamen Wissensgraphen abgebildet werden können. Ziel ist es, klinische Entitäten, Empfehlungen, numerische Zielwerte und deren Beziehungen in eine strukturierte Form zu überführen, die leitlinienübergreifend abgefragt werden kann. Darüber hinaus wird untersucht, ob der resultierende Graph als kontrollierte Grundlage für die Generierung von Lernfragen dienen kann, deren antwortrelevante Inhalte in den extrahierten Informationen verankert sind.

Das vorgeschlagene System besteht aus vier Stufen: Dokumentenaufnahme, Aufbau des Wissensgraphen, Normalisierung und graphgestützte Fragengenerierung. Es wurde anhand von acht Leitliniendokumenten von vier Organisationen evaluiert, die Typ-1- und Typ-2-Diabetes abdecken. Der resultierende Graph umfasst 2.503 Knoten und 3.485 Kanten über 18 Relationstypen. Die Normalisierung von Entitäten und Maßeinheiten ergab acht leitlinienübergreifende Äquivalenzrelationen zwischen ausgewählten numerischen Zielwerten. Der Durchlauf zur Fragengenerierung mit GPT-4o-mini erzeugte 1.289 Fragen, von denen 1.280 die implementierten Validierungsprüfungen bestanden. Der akzeptierte Katalog enthält 3.048 gespeicherte Umformulierungen und 294 Multiple-Choice-Fragen. Eine separate Stichprobe von 30 Fragen aus einem früheren Teildurchlauf mit
Claude Opus 4.8 wurde von zwei Ärztinnen bzw. Ärzten begutachtet; drei eindeutige Fehler ließen sich auf fehlerhafte Relationen im Graphen zurückführen, während eine weitere Frage eher einen Unterschied in der klinischen Praxis als einen sachlichen Fehler aufzeigte.

Die Ergebnisse zeigen, dass der Wissensgraph Informationen aus unterschiedlich strukturierten Diabetesleitlinien zusammenführen und leitlinienübergreifende Abfragen unterstützen kann. Zudem kann der Graph eine kontrolliertere Grundlage für die Generierung von Lernfragen bieten als die direkte Generierung aus dem Leitlinientext. Die Ergebnisse zeigen jedoch auch, dass die Qualität der generierten Fragen von der Qualität des extrahierten Graphen abhängt, insbesondere bei klinischen Relationen. Das Framework ist daher als Grundlage für Anwendungen wie den leitlinienübergreifenden Vergleich, die strukturierte Suche nach klinischen Informationen und die Generierung von Lernfragen gedacht und nicht als Ersatz für die ursprünglichen klinischen Leitlinien.